% =====================================================================
\chapter{El producto de matrices}
\label{tema:4}
% =====================================================================

El producto de matrices se define de modo que represente la composición de aplicaciones lineales. De esa definición se siguen sus propiedades (asociativa, distributivas, comportamiento con la traspuesta) y su principal diferencia con el producto de números: no es conmutativo. El tema sigue la sección~2.4 de Strang~\cite{strang} y la sesión «Multiplication and Inverse Matrices» de 18.06SC~\cite{mit1806sc}; la composición y el coste computacional se tratan como en~\cite[§10.1 y §10.2]{vmls}.

\section{Definición y composición}

\begin{definicion}[Producto de matrices]\label{def:producto}
Sean $A\in\R^{m\times n}$ y $B\in\R^{n\times p}$. El producto $AB\in\R^{m\times p}$ es la matriz de entradas
\[
(AB)_{ij}=\sum_{k=1}^{n}A_{ik}B_{kj},\qquad 1\le i\le m,\ 1\le j\le p.
\]
La entrada $(i,j)$ es el producto escalar de la fila $i$ de $A$ por la columna $j$ de $B$. El producto solo está definido si el número de columnas de $A$ coincide con el número de filas de $B$: $(m\times n)(n\times p)=m\times p$.
\end{definicion}

\begin{teorema}[El producto es la composición]\label{teo:composicion}
Sean $A\in\R^{m\times n}$ y $B\in\R^{n\times p}$. Para todo $x\in\R^{p}$, $(AB)x=A(Bx)$. Es decir, la matriz de la aplicación «aplicar primero $B$ y después $A$» es $AB$.
\end{teorema}

El orden de escritura es el de la composición de funciones, $(A\circ B)(x)=A(B(x))$: en $AB$ actúa primero la matriz de la derecha.

\addtocounter{definicion}{1}%

\section{Cuatro formas de calcular un producto}

La \cref{def:producto} calcula el producto entrada a entrada. Strang subraya otras tres formas equivalentes~\cite[§2.4]{strang}, cada una útil en contextos distintos.

\begin{proposicion}\label{prop:cuatro}
Sean $A\in\R^{m\times n}$ con columnas $a_1,\dots,a_n$ y filas $r_1^{\top},\dots,r_m^{\top}$, y $B\in\R^{n\times p}$ con columnas $b_1,\dots,b_p$ y filas $s_1^{\top},\dots,s_n^{\top}$. Entonces:
\begin{enumerate}
  \item (por columnas) la columna $j$ de $AB$ es $Ab_j$;
  \item (por filas) la fila $i$ de $AB$ es $r_i^{\top}B$;
  \item (columnas por filas) $AB=\sum_{k=1}^{n}a_k\,s_k^{\top}$, suma de $n$ matrices $m\times p$.
\end{enumerate}
\end{proposicion}

Una matriz de la forma $a\,s^{\top}$, con $a\in\R^{m}$ y $s\in\R^{p}$, se llama \emph{producto exterior}; todas sus filas son múltiplos de $s^{\top}$. La forma (iii) descompone un producto de matrices en suma de productos exteriores y es la base de la descomposición en valores singulares y de las aproximaciones de bajo rango del módulo~3. En el \cref{ej:intercambio}, $a_1s_1^{\top}+a_2s_2^{\top}=\begin{pmatrix}0&1\\0&3\end{pmatrix}+\begin{pmatrix}2&0\\4&0\end{pmatrix}=AB$.

\section{Propiedades algebraicas}

\begin{proposicion}\label{prop:algebra}
Siempre que los tamaños permitan formar los productos:
\begin{enumerate}
  \item (asociativa) $(AB)C=A(BC)$;
  \item (distributivas) $A(B+C)=AB+AC$ y $(B+C)A=BA+CA$;
  \item (identidad) $I_mA=A=AI_n$ si $A\in\R^{m\times n}$;
  \item (traspuesta) $(AB)^{\top}=B^{\top}A^{\top}$.
\end{enumerate}
\end{proposicion}

La distributiva por la izquierda y la distributiva por la derecha son propiedades distintas: no puede pasarse de una a otra cambiando el lado (ejercicio~\ref{eje:4-distributiva}), porque el producto no es conmutativo.

\begin{observacion}[No conmutatividad]\label{obs:no-conmuta}
En general $AB\neq BA$, incluso si ambos productos están definidos y son del mismo tamaño. Sean $R=\begin{pmatrix}0&-1\\1&0\end{pmatrix}$ (rotación de $90^\circ$), $S=\begin{pmatrix}2&0\\0&0{,}5\end{pmatrix}$ (escalado) y $x=(3,1)$. Escalar y después rotar da $R(Sx)=R(6;\,0{,}5)=(-0{,}5;\,6)$; rotar y después escalar da $S(Rx)=S(-1,3)=(-2;\,1{,}5)$. El orden de las transformaciones importa.
\end{observacion}

\begin{definicion}[Matriz ortogonal]\label{def:ortogonal}
Una matriz $Q\in\R^{n\times n}$ es \emph{ortogonal} si $Q^{\top}Q=I_n$, es decir, si sus columnas son unitarias y ortogonales dos a dos.
\end{definicion}

\begin{proposicion}\label{prop:ortogonal}
Si $Q$ es ortogonal, entonces $(Qu)\cdot(Qv)=u\cdot v$ para todos $u,v\in\R^{n}$. En consecuencia, $Q$ conserva normas, distancias y ángulos.
\end{proposicion}

Las rotaciones $R_\theta$ son ortogonales (ejercicio~\ref{eje:4-ortogonal}), lo que generaliza la \cref{prop:rotacion}. Las codificaciones posicionales rotatorias (RoPE) del módulo~14 aplican precisamente rotaciones de este tipo a las consultas y las claves.

\section{Coste y matriz de Gram}

\begin{proposicion}[Coste del producto]\label{prop:coste}
Calcular $AB$ con $A\in\R^{m\times n}$ y $B\in\R^{n\times p}$ mediante la \cref{def:producto} requiere $mnp$ multiplicaciones y $mp(n-1)$ sumas.
\end{proposicion}

Aunque $(AB)C=A(BC)$, el coste de los dos órdenes puede ser muy distinto (ejercicio~\ref{eje:4-coste}). Elegir bien el orden de los paréntesis es una optimización habitual en el cálculo con tensores.

\begin{definicion}[Matriz de Gram]\label{def:gram}
Sea $X\in\R^{n\times d}$ la matriz cuyas filas son los vectores $x_1^{\top},\dots,x_n^{\top}$. Su \emph{matriz de Gram} es $G=XX^{\top}\in\R^{n\times n}$, con $G_{ij}=x_i\cdot x_j$.
\end{definicion}

La matriz de Gram es simétrica, porque $G^{\top}=(XX^{\top})^{\top}=XX^{\top}$, y su diagonal contiene las normas al cuadrado, $G_{ii}=\norm{x_i}^{2}$. Un único producto de matrices calcula los $n^{2}$ productos escalares entre $n$ vectores.

\addtocounter{definicion}{1}%

\begin{observacion}
En la atención, las consultas y las claves de los $n$ tokens de una secuencia se apilan como filas de matrices $Q,K\in\R^{n\times d_k}$, y todas las puntuaciones se calculan con un único producto $QK^{\top}\in\R^{n\times n}$, que generaliza la matriz de Gram. Las GPU están diseñadas para ejecutar este tipo de productos con gran rapidez (módulo~17).
\end{observacion}

